\documentclass[10pt,a4paper]{amsart}
\usepackage[margin=19mm]{geometry}
\usepackage{amsmath,amssymb,mathtools}
\usepackage[scr=rsfs,cal=euler]{mathalfa}
\usepackage{xfrac}
\usepackage[unicode,hidelinks,bookmarks=false]{hyperref}
\newcommand{\ie}{i.e.\ }
\newcommand{\cf}{cf.\ }
\newcommand{\resp}{respectively\ }
\newcommand{\Subsection}{\subsection}
\providecommand{\nobreakdash}{\nobreak\hbox{-}}
\setlength{\emergencystretch}{2em}
\multlinegap0pt
\usepackage{amssymb}
\usepackage{cancel}
\usepackage{tikz}
\usepackage{dsfont}
\usepackage{enumerate}
\usepackage{siunitx}
\usepackage{xcolor}
\usepackage[normalem]{ulem}
\newtheorem{lemma}{Lemma}
\newtheorem{corollary}{Corollary}
\newcommand{\Ccalh}{\mathcal{C}_h}
\newcommand{\CcalhK}{\Ccalh(\K)}
\newcommand{\Ceqone}{C_{\text{eq},1}}
\newcommand{\Ceqthree}{C_{\text{eq},3}}
\newcommand{\Ceqtwo}{C_{\text{eq},2}}
\newcommand{\Ctilde}{\widetilde{C}}
\newcommand{\F}{F}
\newcommand{\Fcalh}{\mathcal{F}_h}
\newcommand{\Fi}{{\F_{\Ki}}}
\newcommand{\IKone}{I_{\K,1}}
\newcommand{\IKtwo}{I_{\K,2}}
\newcommand{\Idbf}{\underline{\mathbf{I}}}
\newcommand{\Ione}{I_1}
\newcommand{\Itwo}{I_2}
\newcommand{\K}{K}
\newcommand{\Ki}{{\K_i}}
\newcommand{\Lcalk}{\mathcal{L}_k}
\newcommand{\Lcalkmo}{\mathcal{L}_{k-1}}
\newcommand{\Lcalzero}{\mathcal{L}_0}
\newcommand{\Ltwo}{L^2}
\newcommand{\Norm}[1]{{\left\|{#1} \right\|}}
\newcommand{\Pbb}{\mathbb{P}}
\newcommand{\Pbbk}{\mathbb{P}_k}
\newcommand{\Pbbkmo}{\mathbb{P}_{k-1}}
\newcommand{\Pbbkmt}{\mathbb{P}_{k-2}}
\newcommand{\Pbbz}{\Pbb_0}
\newcommand{\PizF}{\Pi_0^{\F}}
\newcommand{\PizFi}{\Pi_0^{\Fi}}
\newcommand{\Sigmah}{\underline{\boldsymbol{\Sigma}}_h}
\newcommand{\Sigmahkmo}{\Sigmah^{k-1}}
\newcommand{\Sigmahz}{\Sigmah^0}
\newcommand{\Tauh}{\mathcal{T}_h}
\newcommand{\Vbf}{\mathbf{V}}
\newcommand{\Vh}{{\bf{V}}_h}
\newcommand{\Vhk}{\Vh^k}
\newcommand{\Vhkz}{\Vbf_{h,0}^k}
\newcommand{\btwo}{b_2}
\newcommand{\ceqone}{c_{\text{eq},1}}
\DeclareMathOperator{\dev}{\underline{\mathbf{dev}}}
\DeclareMathOperator{\dive}{div}
\newcommand{\ds}{\, \mathrm{d}s}
\newcommand{\dx}{\, \mathrm{d}x}
\newcommand{\epsbfun}{\underline{\boldsymbol{\varepsilon}}}
\newcommand{\hKi}{h_{\Ki}}
\newcommand{\itt}{\mathtt{i}}
\newcommand{\jtt}{\mathtt{j}}
\newcommand{\jump}[1]{\left[\!\left[#1\right]\!\right]}
\newcommand{\lambdai}{\lambda_i}
\newcommand{\lambdaop}{\lambda^{\text{op}}}
\newcommand{\lkmo}{\ell_{k-1}}
\newcommand{\mkmo}{m_{k-1}}
\newcommand{\ncompo}{n_1}
\newcommand{\ncompt}{n_2}
\DeclareMathOperator{\norm}{\bf{n}}
\newcommand{\normi}{\norm_i}
\newcommand{\normitt}{\norm_{\itt}}
\newcommand{\normj}{\norm_j}
\newcommand{\normjtt}{\norm_{\jtt}}
\newcommand{\normtang}{{nt}}
\newcommand{\piF}{p_{i}^F}
\newcommand{\psiEi}{\psi^{\Fi}}
\newcommand{\psiKio}{\psi^{\Ki}_1}
\newcommand{\psiKit}{\psi^{\Ki}_2}
\newcommand{\qi}{q_{i}}
\newcommand{\qk}{q_k}
\newcommand{\qkF}{q^\F_{k}}
\newcommand{\qkmo}{q_{k-1}}
\newcommand{\rkmoi}{r_{k-1}^i}
\newcommand{\skmoi}{s_{k-1}^i}
\newcommand{\stildekmt}{\widetilde{s}_{k-2}}
\newcommand{\stildekmti}{\stildekmt^i}
\newcommand{\sym}{\underline{\bf{sym}}}
\DeclareMathOperator{\tang}{\bf{t}}
\newcommand{\tangi}{\tang_i}
\newcommand{\tangj}{\tang_j}
\newcommand{\tangscalar}{t}
\newcommand{\taubf}{\underline{\boldsymbol{\tau}}}
\newcommand{\taubfh}{\underline{\boldsymbol{\tau}}_h}
\DeclareMathOperator{\tr}{tr}
\newcommand{\ubf}{\mathbf{u}}
\newcommand{\vbfh}{\mathbf{v}_h}
\newcommand{\wi}{w_i}
\newcommand{\xbf}{\mathbf{x}}
\title{A mass-conserving stress-yielding formulation for the Stokes equation with exact stress symmetry}
\author{Philip L. Lederer, Marialetizia Mosconi}
\date{Source: arXiv:2609.18382v1 (CC BY 4.0)}
\begin{document}
\maketitle

\noindent\emph{Source and licence.} A mass-conserving stress-yielding formulation for the Stokes equation with exact stress symmetry. Version arXiv:2609.18382v1 (CC BY 4.0), distributed by arXiv under the Creative Commons Attribution 4.0 International licence, http://creativecommons.org/licenses/by/4.0/, as recorded in the arXiv licence metadata for this exact version and shown on the abstract page https://arxiv.org/abs/2609.18382v1. Full licence notice: \texttt{licenses/arxiv-2609.18382-CC-BY-4.0.txt}.\par\smallskip

\noindent\textit{Editorial scope.} Counted unit: the article's lemma lemma:discrete-inf-sup-bt (source0.tex lines 2309--2328). The article's complete original proof of it is delivered with it (source0.tex lines 2329--2767, one \texttt{proof} environment of 439 lines). No mathematics is written, repaired or re-derived: the statement and the proof are the article's own lines, in the article's own notation, and no retained line is substituted or re-flowed.  Retained uncounted as the source-local context the proof needs: lines 1049--1051; lines 1057--1060; lines 1224--1228; lines 1355--1362; lines 1364--1371; lines 1387--1396; lines 1417--1429; lines 1445--1483; lines 1526--1528; lines 1793--1813; lines 1895--1910; lines 1920--1925; lines 2144--2146; lines 2155--2167; lines 2220--2224.  Nothing was omitted from any retained span, so the package carries no omission record.  No retained line contains a citation, so the wrapper carries no bibliography and no bibliography entry.  No retained line contains a figure environment, an image-inclusion command or drawing code, so no image is delivered and the package declares no asset dependency.  Editorial formatting only, in addition: an AMS \texttt{amsart} preamble that restates the article's own declarations and macros used by the retained lines, namely the environments \texttt{lemma}, \texttt{corollary} on separate counters, together with the standard packages those lines need.  The retained environments print in this document as Lemma 1, Corollary 1 in order of appearance, whereas the article prints the same items with the numbering of its own full document.  The complete native source of the article as distributed in the arXiv e-print for this exact version is bundled unchanged as \texttt{sources/210696/source0.tex}.
\begin{equation} \label{Lcal}
\Lcalk(\qkF)(\xbf) := \qkF(\xbf_{\Fi}),
\end{equation}

\begin{equation} \label{stability-Lcal}
\Norm{\Lcalk(\qkF)}_{\Ltwo(\Ki)}
\lesssim \hKi^{\frac12} \Norm{\qkF}_{\Ltwo(\Fi)}.
\end{equation}

\begin{equation} \label{btwo}
\btwo(\taubf,\ubf)
:= - \sum_{\K \in \Tauh} \sum_{i=1}^3 \int_{\Ki} \taubf : \epsbfun(\ubf)\dx 
+ \sum_{\F \in \Fcalh} \int_{\F} \taubf_{\normtang} \jump{\ubf_t} \ds.
\end{equation}

\begin{equation} \label{dev:frob-norm}
\begin{split}
|\dev(\normi\otimes\normi)|^2
&= \dev(\normi\otimes\normi):\dev(\normi\otimes\normi) \\
&=  (\ncompo^2 - \frac12)^2 + 2\ncompo^2\ncompt^2 + (\ncompt^2 - \frac12)^2
=  (\ncompo^2 + \ncompt^2)^2 - \ncompo^2 - \ncompt^2 + \frac12 = \frac12 
\end{split}
\end{equation}

\begin{equation} \label{sym:frob-norm}
\begin{split}
|\sym(\normi\otimes\tangi)|^2
&=\sym(\normi\otimes\tangi):\sym(\normi\otimes\tangi) \\
&=  2\ncompo^2\ncompt^2 + 2 \big(\frac12 (\ncompt^2 -\ncompo^2)\big)^2   
=  \frac12\big(\ncompo^2 + \ncompt^2\big)^2 = \frac12.
\end{split}
\end{equation}

\begin{equation} \label{orthogonality:dev-sym}
\begin{split}
\dev(\normi\otimes\normi) : \sym(\normi\otimes\tangi)
&= \dev(\normi\otimes\normi) : (\normi\otimes\tangi) \\
&= (\normi\otimes\normi) : (\normi\otimes\tangi)
    - \frac{1}{2}\tr(\normi\otimes\normi)\Idbf:(\normi\otimes\tangi) \\
&= (\normi\cdot\normi)(\normi\cdot \tangi) - \frac{1}{2}(\normi\cdot \tangi)
= 0.
\end{split}
\end{equation}

\begin{lemma} \label{lemma:dev-sym-basis-Ki}
Let~$\Ki$ be an element of~$\Ccalh$.
Then,
the set $\{\dev(\normi\otimes\normi),\sym(\normi\otimes\tangi)\}$
is a basis for~$\Sigmahz(\Ki)$.
In particular,
for any given~$\taubfh$ in~$\Sigmahkmo(\Ki)$,
there exist polynomials~$\rkmoi$ and~$\skmoi$ in~$\Pbbkmo(\Ki)$
such that~$\taubfh$ admits the following representation:
\begin{equation*}
\taubfh = \rkmoi\dev(\normi\otimes\normi) + \skmoi\sym(\normi\otimes\tangi).
\end{equation*}
\end{lemma}

We now aim at introducing a set of degrees of freedom
for~$\Sigmahkmo(\Ki)$.
For a given~$\Ki$ in~$\Ccalh$
with external edge~$\Fi$,
consider the following three groups of linear functionals:
\begin{subequations} \label{dofs}
\begin{itemize}
\item
edge moments~$\psiEi(\taubfh)$ on~$\Fi$
associated with the normal-tangential component,
up to order~$k-1$:
    \begin{equation} \label{dofs:edge}
    \psiEi(\taubfh):= \Big\{
            \int_{\Fi} (\taubfh)_{\normtang}  \piF \ds,\,
            \forall \piF \in \Pbbkmo(\Fi)
                    \Big\};
    \end{equation}
\item bulk moments~$\psiKio(\taubfh)$
associated with the deviatoric component,
up to order~$k-1$:
    \begin{equation} \label{dofs:bulk-k}
    % \psiKo(\taubfh) := \Big\{ \int_\F \taubfh : \dev(\norm \otimes \norm)\qk \dx,\,
    % \forall\qk \in \Pbbk(\K) \Big\}
    \psiKio(\taubfh) := \Big\{ \int_{\Ki} \taubfh : \dev(\normi\otimes\normi) \qi \dx,\,
    \forall\qi \in \Pbbkmo(\Ki) \Big\};
    \end{equation}
\item bulk moments~$\psiKit(\taubfh)$
associated with the normal-tangential component weighted by the barycenter coordinate,
up to order~$k-2$:
    \begin{equation} \label{dofs:bulk-kmo-lambda}
    % \psiKt(\taubfh) := \Big\{ \int_\F \taubfh : (\norm \otimes \tang) \qkmo \lambdaop \dx,\,
    % \forall\qkmo \in \Pbbkmo(\K) \Big\}
    \psiKit(\taubfh) := \Big\{ \int_{\Ki} \taubfh :
    \sym(\normi\otimes\tangi)
    \wi \lambdai \dx,\,
    \forall\wi \in \Pbbkmt(\Ki) \Big\};
\end{equation}
\end{itemize}
\end{subequations}

\begin{equation} \label{decomposition-polynomial-spaces}
\Pbbkmo(\Ki) = {\Lcalkmo(\Pbbkmo(\Fi))}\oplus\lambdai\Pbbkmt(\Ki).
\end{equation}

\begin{lemma} \label{lemma:dev-n1-n2-basis}
Given~$\K$ in~$\Tauh$,
there exist two distinct indices $\itt,\jtt$ in $\{1,2,3\}$
such that the set
$\{\dev(\normitt\otimes\normitt),\dev(\normjtt\otimes\normjtt)\}$
forms a basis for~$\Sigmahz(\K)$.
In particular,
for any given~$\taubfh$ in~$\Sigmahkmo(\K)$,
there exist polynomials~$\lkmo$ and~$\mkmo$ in~$\Pbbkmo(\K)$
such that~$\taubfh$ admits the following representation:
\begin{equation*}
\taubfh = \lkmo \dev(\normitt\otimes\normitt) + \mkmo \dev(\normjtt\otimes\normjtt).
\end{equation*}
Moreover,
the following estimate holds true
\begin{equation} \label{strict-CS}
|\dev(\normitt\otimes\normitt):\dev(\normjtt\otimes\normjtt)|
< |\dev(\normitt\otimes\normitt)||\dev(\normjtt\otimes\normjtt)|
= \frac12.
\end{equation}
\end{lemma}

\begin{corollary} \label{corollary:decomposition-epsvh}
Given~$\K$ in~$\Tauh$
and~$\vbfh$ in~$\Vhk$,
there exist polynomials~$\lkmo$ and~$\mkmo$ in~$\Pbbkmo(\K)$
such that the restriction of~$\dev(\epsbfun(\vbfh))$ to~$\K$
admits the following representation:
\begin{equation} \label{eps-vh:dev-dev-representation-dev}
\dev(\epsbfun(\vbfh)|_{\K}) = \lkmo \dev(\normitt\otimes\normitt) + \mkmo \dev(\normjtt\otimes\normjtt).
\end{equation}
Moreover,
if~$\dive(\vbfh|_{\K})$ vanishes,
the same representation holds also for the restriction of~$\epsbfun(\vbfh)$ to~$\K$, i.e.,
\begin{equation} \label{eps-vh:dev-dev-representation}
\epsbfun(\vbfh)|_{\K} = \lkmo \dev(\normitt\otimes\normitt) + \mkmo \dev(\normjtt\otimes\normjtt).
\end{equation}
\end{corollary}

\begin{equation} \label{deveps-equal-eps}
\dev(\epsbfun(\vbfh))
= \epsbfun(\vbfh) - \frac{1}{d}\tr(\epsbfun(\vbfh)) \Idbf 
= \epsbfun(\vbfh) - \frac{1}{d}\dive(\vbfh) \Idbf 
= \epsbfun(\vbfh).
\end{equation}

\begin{equation} \label{norm-equivalence-h-dev}
\Norm{\vbfh}_{\Vh} \simeq \Norm{\vbfh}_{1,\dev,h},
\end{equation}

\begin{lemma} \label{lemma:norm-equivalence}
Let~$\K$ be an element in~$\Tauh$
and~$\vbfh$ be in~$\Vhkz$.
Consider its symmetric gradient~$\epsbfun(\vbfh)$
and let~$\lkmo$ and~$\mkmo$ be the polynomials in~$\Pbbkmo(\K)$
of its representation on~$\K$; cf.~\eqref{eps-vh:dev-dev-representation-dev}.
Then,
the following norm equivalence holds true
\begin{equation} \label{norm-equivalence:statement}
\Norm{\epsbfun(\vbfh)}_{\Ltwo(\K)}^2
\simeq \sum_{i = 1}^3\big(\Norm{\lkmo}_{\Ltwo(\Ki)}^2 + \Norm{\mkmo}_{\Ltwo(\Ki)}^2\big).
\end{equation}
\end{lemma}

\begin{equation} \label{norm-equivalence:explicit-constants}
\ceqone\Norm{\epsbfun(\vbfh)}_{\Ltwo(\K)}
\leq \Big(\sum_{i = 1}^3\big(\Norm{\lkmo}_{\Ltwo(\Ki)}^2 + \Norm{\mkmo}_{\Ltwo(\Ki)}^2\big)\Big)^{\frac12}
\leq \Ceqone\Norm{\epsbfun(\vbfh)}_{\Ltwo(\K)}.
\end{equation}

\begin{lemma} \label{lemma:discrete-inf-sup-bt}
For any~$\vbfh$ in~$\Vhkz$,
there exists~$\taubfh$ in~$\Sigmahkmo$ such that
\begin{subequations} \label{inf-sup:i-ii}
\begin{equation} \label{inf-sup:i}
\btwo(\taubfh,\vbfh) \gtrsim \Norm{\vbfh}_{1,\dev,h}^2
\end{equation}
\begin{equation} \label{inf-sup:ii}
\Norm{\taubfh}_{\Sigmah} \lesssim \Norm{\vbfh}_{1,\dev,h}.
\end{equation}
\end{subequations}
These estimates imply
the discrete inf-sup stability of the bilinear form~$\btwo(\cdot,\cdot)$ of~\eqref{btwo} on~$\Vhkz$,
i.e.,
for all~$\vbfh$ in $\Vhkz$,
\begin{equation} \label{discrete-inf-sup:btwo}
\sup_{\taubfh \in \Sigmahkmo} \frac{\btwo(\taubfh,\vbfh)}{\Norm{\taubfh}_{\Sigmah}}
    \gtrsim\Norm{\vbfh}_{1,h}.
\end{equation}
\end{lemma}

\begin{proof}
We detail the proof for $k-1 \geq 1$.
The lowest-order case $k-1 = 0$ can be treated similarly
and is simpler,
since all the estimates below that rely on the combination of
the degrees of freedom~\eqref{dofs:edge} and~\eqref{dofs:bulk-kmo-lambda}
can be directly obtained by using just the edge moments~\eqref{dofs:edge}.

Let~$\vbfh$ be any given function in~$\Vhkz$;
we proceed by constructing a stress function~$\taubfh$ in~$\Sigmahkmo$
that satisfies the estimates~\eqref{inf-sup:i}--\eqref{inf-sup:ii}.
Let~$\gamma > 0$ be a positive constant to be fixed later.
We define~$\taubfh$
by fixing the corresponding degrees of freedom in~\eqref{dofs} as follows.
We impose:
\begin{itemize}
\item on every edge~$\F$ in~$\Fcalh$, for all~$\piF$ in~$\Pbbkmo(\F)$,
\begin{equation}  \label{stability-dofs:1}
\int_{\F} (\taubfh)_{\normtang} \piF \ds = \frac{1}{h} \int_{\F} \PizF\jump{(\vbfh)_{\tangscalar}}\piF \ds;
\end{equation}
\item on every subelement~$\Ki$ in~$\Ccalh$, for all~$\qi$ in~$\Pbbkmo(\Ki)$,
\begin{equation} \label{stability-dofs:2}
\int_{\Ki} \taubfh : \dev(\normi\otimes\normi) \qi \dx = - \gamma\int_{\Ki} \epsbfun(\vbfh) : \dev(\normi\otimes\normi) \qi \dx;
\end{equation}
\item on every subelement~$\Ki$ in~$\Ccalh$, for all~$\wi$ in~$\Pbbkmt(\Ki)$,
\begin{equation} \label{stability-dofs:3}
\int_{\Ki} \taubfh : \sym(\normi\otimes\tangi)\wi\lambdai \dx = 0.
\end{equation}
\end{itemize}
Now let~$\K$ be any fixed element in~$\Tauh$,
and let~$\Ki$ be in~$\CcalhK$.
Since~$\taubfh$ belongs to~$\Sigmahkmo$,
as discussed in Lemma~\ref{lemma:dev-sym-basis-Ki},
there exist polynomials~$\rkmoi$ and~$\skmoi$ in~$\Pbbkmo(\Ki)$
such that the restriction of~$\taubfh$ on~$\Ki$
admits the following representation:
\begin{equation} \label{tau-decomposition}
\taubfh|_{\Ki} = \rkmoi\dev(\normi\otimes\normi) + \skmoi\sym(\normi\otimes\tangi).
\end{equation} 
Similarly,
since~$\epsbfun(\vbfh)$ belongs to~$\Sigmahkmo(\K)$,
as discussed in Corollary~\ref{corollary:decomposition-epsvh},
there exist polynomials~$\lkmo$ and $\mkmo$ in~$\Pbbkmo(\K)$
such that the restriction of~$\epsbfun(\vbfh)$ on~$\K$
admits the following representation:
\begin{equation} \label{eps-decomposition}
\epsbfun(\vbfh)|_{\K} = \lkmo\dev(\normitt\otimes\normitt) + \mkmo\dev(\normjtt\otimes\normjtt).
\end{equation}
Since the subelements~$\Ki \in \CcalhK$ are subsets of $\K$,
we can restrict the representation~\eqref{eps-decomposition}
to each~$\Ki$.

We now give alternative representations
for the coefficients of the decomposition~\eqref{tau-decomposition} above,
starting from~$\skmoi$.
Since $(\taubfh)_{\normtang}|_{\F}$
belongs to~$\Pbbkmo(\F)$,
identity~\eqref{stability-dofs:1},
entails that,
for all~$\F$ in~$\Fcalh$,
\begin{equation*}
\begin{split}
\frac{1}{h}\PizF\jump{(\vbfh)_{\tangscalar}}
\overset{\eqref{stability-dofs:1}}{=}
(\taubfh)_{\normtang}|_{\F} 
&\overset{\eqref{tau-decomposition}}{=}
\big(\rkmoi\dev(\normi\otimes\normi) + \skmoi\sym(\normi\otimes\tangi)\big)_{\normtang}|_{\F} 
\overset{\eqref{orthogonality:dev-sym},\eqref{sym:frob-norm}}{=}
\frac{1}{2}\skmoi|_{\F}.
\end{split}
\end{equation*}
In particular,
for each~$\Ki$ in~$\CcalhK$
with corresponding external edge~$\Fi$ in~$\Fcalh$,
it holds
\begin{equation}  \label{skmoi-value-on-edge}
\frac{1}{h}\PizFi\jump{(\vbfh)_{\tangscalar}}
= \frac{1}{2}\skmoi|_{\Fi}.
\end{equation}
Since~$\skmoi$ belongs to~$\Pbbkmo(\Ki)$,
the above display combined with the decomposition~\eqref{decomposition-polynomial-spaces},
yield the existence of a polynomial~$\stildekmti$ in~$\Pbbkmt(\Ki)$
such that
\begin{equation} \label{skmoi-stildekmti}
\skmoi = \frac{2}{h}\Lcalzero(\PizFi\jump{(\vbfh)_{\tangscalar}}) + \stildekmti\lambdai,
\end{equation}
where
$\Lcalzero(\PizFi\jump{(\vbfh)_{\tangscalar}})$ denotes
the constant extension of $\PizFi\jump{(\vbfh)_{\tangscalar}}$ from~$\Fi$ over~$\Ki$
as defined in~\eqref{Lcal}.

We now also give a representation for the polynomial coefficient~$\rkmoi$ of~\eqref{tau-decomposition}.
Inserting the \eqref{tau-decomposition}
in~\eqref{stability-dofs:2},
and applying the orthogonality property~\eqref{orthogonality:dev-sym},
we obtain that,
for all~$\qi$ in $\Pbbkmo(\Ki)$,
the following identity holds true
\begin{equation*} \label{rkmoi-integral-value}
\int_{\Ki} \rkmoi \dev(\normi\otimes\normi) : \dev(\normi\otimes\normi) \qi \dx = - \gamma\int_{\Ki} \epsbfun(\vbfh) : \dev(\normi\otimes\normi) \qi \dx.
\end{equation*}
Since the polynomial coefficient $\rkmoi$
belongs to~$\Pbbkmo(\Ki)$
and the symmetric gradient $\epsbfun(\vbfh)|_{\Ki}$
and the tensor $\dev(\normi\otimes\normi)$
have entries in~$\Pbbkmo(\Ki)$, and~$\Pbbz(\Ki)$,
respectively,
the above display
combined with the identity~\eqref{dev:frob-norm}
yields
\begin{equation} \label{rkmoi-value}
\begin{split}
\rkmoi
&= -2\gamma \dev(\normi\otimes\normi) : \epsbfun(\vbfh)|_{\Ki}.
\end{split}
\end{equation}

We are now ready to prove the estimates in~\eqref{inf-sup:i-ii},
starting with~\eqref{inf-sup:i}.
Plugging~$\taubfh$ and~$\vbfh$
in the bilinear form~$\btwo(\cdot,\cdot)$ of~\eqref{btwo},
we have
\begin{equation*} \label{bt-Ione-Itwo}
\begin{split}
\btwo(\taubfh,\vbfh) 
= -\sum_{\K \in \Tauh} \sum_{i=1}^3 \int_{\Ki} \taubfh : \epsbfun(\vbfh) \dx
    + \sum_{\F \in \Fcalh} \int_{\F} (\taubfh)_{\normtang}\jump{(\vbfh)_{\tangscalar}} \ds 
=: \Ione + \Itwo.
\end{split}
\end{equation*}

We start by estimating~$\Ione$.
Consider the constants
$B_{ij} := \dev(\normi\otimes\normi):\sym(\normj\otimes\tangj)$;
in particular,
Cauchy--Schwarz inequality gives~$|B_{ij}| \leq 1$.
Note that below,
one of the two indices will always be fixed to~$\itt$ or~$\jtt$
of Lemma~\ref{lemma:dev-n1-n2-basis}.
Using the representations~\eqref{tau-decomposition} of~$\taubfh$
and~\eqref{eps-decomposition} of~$\epsbfun(\vbfh)$  on each~$\Ki$,
the orthogonality~\eqref{orthogonality:dev-sym},
and the identity~\eqref{rkmoi-value},
we have
\begin{equation} \label{Ione}
\begin{split}
-\sum_{i=1}^3& \int_{\Ki} \taubfh : \epsbfun(\vbfh) \dx 
\overset{\eqref{tau-decomposition}}{=}
-\sum_{i=1}^3 \int_{\Ki} (\rkmoi\dev(\normi\otimes\normi)+\skmoi\sym(\normi\otimes\tangi)) : \epsbfun(\vbfh) \dx \\
&\overset{\eqref{orthogonality:dev-sym},\eqref{eps-decomposition}}{=} 
-\sum_{i=1}^3 \int_{\Ki} 
   \rkmoi\,(\dev(\normi\otimes\normi):\epsbfun(\vbfh))
   + \skmoi\,(\lkmo B_{i\itt} + \mkmo B_{i\jtt}) \dx \\
&\overset{\eqref{rkmoi-value}}{=}
\sum_{i=1}^3 \int_{\Ki} 
   2\gamma(\dev(\normi\otimes\normi):\epsbfun(\vbfh))^2 
   \dx
   -\sum_{i=1}^3 \int_{\Ki} \skmoi\,(\lkmo B_{i\itt} + \mkmo B_{i\jtt}) \dx\\
&=: {\IKone} + {\IKtwo}.
\end{split}
\end{equation}
We estimate the two terms on the right-hand side separately.

We first estimate~$\IKone$.
Since~$\dev(\normi\otimes\normi):\epsbfun(\vbfh)$ is a polynomial,
a scaling argument combined with~$|\Ki|=|K|/3$
and the shape-regularity assumption
yield
that there exists a positive constant~$\Ceqtwo$ such that
\begin{equation*}
\Norm{\dev(\normi\otimes\normi):\epsbfun(\vbfh)}_{\Ltwo(\K)}
\leq \Ceqtwo \Norm{\dev(\normi\otimes\normi):\epsbfun(\vbfh)}_{\Ltwo(\Ki)},
\end{equation*}
which implies
\begin{equation} \label{IKone:first-estimate}
\IKone 
= \sum_{i=1}^3 2\gamma
  \Norm{\dev(\normi\otimes\normi):\epsbfun(\vbfh)}_{\Ltwo(\Ki)}^2
\geq 2\gamma\Ceqtwo^2 \sum_{i=1}^3  
  \Norm{\dev(\normi\otimes\normi):\epsbfun(\vbfh)}_{\Ltwo(\K)}^2.
\end{equation}
Consider the constants
$D_{ij} := \dev(\normi\otimes\normi):\dev(\normj\otimes\normj)$;
with this notation,
identity~\eqref{dev:frob-norm} reads as $|D_{ii}|=1/2$.
Note that below,
the two indices will always be fixed to~$\itt$ and~$\jtt$ of Lemma~\ref{lemma:dev-n1-n2-basis}.
Neglecting a non-negative term,
using the representation~\eqref{eps-decomposition} of~$\epsbfun(\vbfh)$ on~$\K$
and the definition of~$D_{ij}$,
and expanding the squares,
we obtain
\begin{equation} \label{sumdev-eps}
\begin{split}
\sum_{i = 1}^3&\Norm{\dev(\normi\otimes\normi):\epsbfun(\vbfh)}_{\Ltwo(\K)}^2 
\geq \sum_{i \in\itt,\jtt}\Norm{\dev(\normi\otimes\normi):\epsbfun(\vbfh)}_{\Ltwo(\K)}^2 \\
&\overset{\eqref{eps-decomposition}}{=} 
\sum_{i \in\itt,\jtt}\Norm{\dev(\normi\otimes\normi):(\lkmo\dev(\normitt\otimes\normitt)+\mkmo\dev(\normjtt\otimes\normjtt))}_{\Ltwo(\K)}^2 \\
&=\Norm{\lkmo D_{\itt\itt} + \mkmo D_{\itt\jtt}}_{\Ltwo(\K)}^2 
     + \Norm{\lkmo D_{\itt\jtt}+ \mkmo D_{\jtt\jtt}}_{\Ltwo(\K)}^2 \\
&=\Norm{\frac12\lkmo + \mkmo D_{\itt\jtt}}_{\Ltwo(\K)}^2 
     + \Norm{\lkmo D_{\itt\jtt}+ \frac12\mkmo}_{\Ltwo(\K)}^2 \\
&= \frac14\Norm{\lkmo}^2_{\Ltwo(\K)} + |D_{\itt\jtt}|^2\Norm{\mkmo}^2_{\Ltwo(\K)} + \int_\K\lkmo\mkmo D_{\itt\jtt}\dx \\
        &\quad +|D_{\itt\jtt}|^2\Norm{\lkmo}^2_{\Ltwo(\K)} + \frac14\Norm{\mkmo}^2_{\Ltwo(\K)} +\int_\K\lkmo\mkmo D_{\itt\jtt}\dx \\
&=\Big(\frac14+|D_{\itt\jtt}|^2\Big)(\Norm{\lkmo}^2_{\Ltwo(\K)} + \Norm{\mkmo}^2_{\Ltwo(\K)})+ 2\int_\K\lkmo\mkmo D_{\itt\jtt}\dx \\
\end{split}
\end{equation}
We estimate the last term on the right-hand side.
The Cauchy--Schwarz inequality and
the Young inequality
yield
\begin{equation*}
\begin{split}
2\int_\K\lkmo\mkmo D_{\itt\jtt}\dx
&\geq - 2\int_\K|\lkmo\mkmo D_{\itt\jtt}|\dx 
\geq - 2|D_{\itt\jtt}|\Norm{\lkmo}_{\Ltwo(\K)}\Norm{\mkmo}_{\Ltwo(\K)} \\
&\geq - 2|D_{\itt\jtt}|\big(\frac12\Norm{\lkmo}^2_{\Ltwo(\K)} + \frac12\Norm{\mkmo}^2_{\Ltwo(\K)}\big).
\end{split}
\end{equation*}
Using the estimate above in~\eqref{sumdev-eps}
and applying Lemma~\ref{lemma:norm-equivalence}
with the notation of~\eqref{norm-equivalence:explicit-constants},
gives
\begin{equation*}
\begin{split}
&\sum_{i = 1}^3\Norm{\dev(\normi\otimes\normi):\epsbfun(\vbfh)}_{\Ltwo(\K)}^2 
\geq \Big(\frac14+|D_{\itt\jtt}|^2-|D_{\itt\jtt}|\Big)(\Norm{\lkmo}^2_{\Ltwo(\K)} + \Norm{\mkmo}^2_{\Ltwo(\K)}) \\
&\qquad= \Big(\frac12-|D_{\itt\jtt}|\Big)^2(\Norm{\lkmo}^2_{\Ltwo(\K)} + \Norm{\mkmo}^2_{\Ltwo(\K)})
\geq \Big(\frac12-|D_{\itt\jtt}|\Big)^2\ceqone^2\Norm{\epsbfun(\vbfh)}^2_{\Ltwo(\K)} \\
&\qquad=: \Ceqthree^2\Norm{\epsbfun(\vbfh)}^2_{\Ltwo(\K)}.
\end{split}
\end{equation*}
Note that due to~\eqref{strict-CS}, $|D_{\itt\jtt}|<1/2$ and
consequently, $\Ceqthree$ is a positive constant.
Inserting the above display in~\eqref{IKone:first-estimate},
we obtain
\begin{equation} \label{IKone}
\IKone \geq 2\gamma\Ceqtwo^2\Ceqthree^2\Norm{\epsbfun(\vbfh)}_{\Ltwo(\K)}^2.
\end{equation}

We now estimate~$\IKtwo$.
Inserting the representations~\eqref{tau-decomposition} of~$\taubfh$
and~\eqref{skmoi-stildekmti} of~$\skmoi$ on~$\Ki$
in identity~\eqref{stability-dofs:3},
we obtain that,
for all~$\wi$ in~$\Pbbkmt(\Ki)$
it holds
\begin{equation} \label{zero-sum-Lcal-skmoi}
\begin{split}
0
&\overset{\eqref{stability-dofs:3}}{=}
\int_{\Ki} \taubfh : \sym(\normi\otimes\tangi)\wi\lambdai \dx \\
&\overset{\eqref{tau-decomposition}}{=}
\int_{\Ki} (\rkmoi\dev(\normi\otimes\normi) + \skmoi\sym(\normi\otimes\tangi)) : \sym(\normi\otimes\tangi)\wi\lambdai \dx \\
&\overset{\eqref{orthogonality:dev-sym},\eqref{skmoi-stildekmti}}{=}
\int_{\Ki} 0 + \big(\frac{2}{h}\Lcalzero(\PizFi\jump{(\vbfh)_{\tangscalar}}) + \stildekmti\lambdai\big) \sym(\normi\otimes\tangi) : \sym(\normi\otimes\tangi)\wi\lambdai \dx \\
&\overset{\eqref{sym:frob-norm}}{=}
\frac12 \int_{\Ki} \big(\frac{2}{h}\Lcalzero(\PizFi\jump{(\vbfh)_{\tangscalar}}) + \stildekmti\lambdai\big)\wi\lambdai \dx.
\end{split}
\end{equation}
Choosing
$\wi = \stildekmti$ in the above display
and applying
the Cauchy--Schwarz inequality
and the stability estimate~\eqref{stability-Lcal} of the constant extension operator~$\Lcalzero$,
give
\begin{equation*}
\begin{split}
\Norm{\stildekmti\lambdai}_{\Ltwo(\Ki)}^2
&\overset{\eqref{zero-sum-Lcal-skmoi}}{\lesssim}
\frac{1}{h}\int_{\Ki}|\Lcalzero(\PizFi\jump{(\vbfh)_{\tangscalar}})\stildekmti\lambdai| \dx \\
&\leq \frac{1}{h}\Norm{\Lcalzero(\PizFi\jump{(\vbfh)_{\tangscalar}})}_{\Ltwo(\Ki)}\Norm{\stildekmti\lambdai}_{\Ltwo(\Ki)} \\
&\overset{\eqref{stability-Lcal}}{\lesssim}
h^{-\frac12}\Norm{\PizFi\jump{(\vbfh)_{\tangscalar}}}_{\Ltwo(\Fi)}\Norm{\stildekmti\lambdai}_{\Ltwo(\Ki)}.
\end{split} 
\end{equation*}
Then,
taking the $\Ltwo$ norm on both sides of the identity~\eqref{skmoi-stildekmti}
and applying 
the triangle inequality,
the above display,
the stability estimate~\eqref{stability-Lcal},
and the shape-regularity assumption,
yield
the existence of a positive constant~$\Ctilde$
such that
\begin{equation} \label{skmoi-leq-PizF}
\begin{split}
\Norm{\skmoi}_{\Ltwo(\Ki)}
&\leq\Norm{\stildekmti\lambdai}_{\Ltwo(\Ki)} + \frac{2}{h}\Norm{\Lcalzero(\PizFi\jump{(\vbfh)_{\tangscalar}})}_{\Ltwo(\Ki)} \\
&\leq \Ctilde h^{-\frac12}\Norm{\PizFi\jump{(\vbfh)_{\tangscalar}}}_{\Ltwo(\Fi)}.
\end{split}
\end{equation}
Considering now the term~$\IKtwo$ in~\eqref{Ione},
applying the Cauchy--Schwarz inequality,
the triangle inequality,
the estimates~\eqref{skmoi-leq-PizF}
and~$|B_{ij}| \leq 1$ for~$ij = \{i\itt,i\jtt\}$,
the weighted Young inequality with parameter~$\delta > 0$,
and Lemma~\ref{lemma:norm-equivalence}
with the notation of~\eqref{norm-equivalence:explicit-constants},
we have
\begin{equation} \label{IKtwo}
\begin{split}
|\IKtwo|
&\leq \sum_{i=1}^3 \Norm{\skmoi}_{\Ltwo(\Ki)} \Norm{\lkmo B_{i\itt} + \mkmo B_{i\jtt}}_{\Ltwo(\Ki)} \\
&\leq \sum_{i=1}^3 \Norm{\skmoi}_{\Ltwo(\Ki)} \big(\Norm{\lkmo B_{i\itt}}_{\Ltwo(\Ki)} + \Norm{\mkmo B_{i\jtt}}_{\Ltwo(\Ki)}\big) \\
&\leq \sum_{i=1}^3 \Ctilde \frac{1}{h^{\frac12}} \Norm{\PizFi\jump{(\vbfh)_{\tangscalar}}}_{\Ltwo(\Fi)} \big(\Norm{\lkmo}_{\Ltwo(\Ki)} + \Norm{\mkmo}_{\Ltwo(\Ki)}\big) \\
&\leq \Ctilde\frac{1}{2\delta}\sum_{i=1}^3 \frac{1}{h} \Norm{\PizFi\jump{(\vbfh)_{\tangscalar}}}^2_{\Ltwo(\Fi)}
        + \delta\sum_{i=1}^3 \big(\Norm{\lkmo}_{\Ltwo(\Ki)}^2 + \Norm{\mkmo}_{\Ltwo(\Ki)}^2\big) \\
&= \frac{\Ctilde}{2\delta}\sum_{i=1}^3 \frac{1}{h} \Norm{\PizFi\jump{(\vbfh)_{\tangscalar}}}^2_{\Ltwo(\Fi)}
        + \delta\Ceqone^2\Norm{\epsbfun(\vbfh)}_{\Ltwo(\K)}^2.
\end{split}
\end{equation}
Summing over the elements~$\K$ in~$\Tauh$
and plugging~\eqref{IKone} and~\eqref{IKtwo} in~\eqref{Ione},
since one edge in~$\Fcalh$ is shared at most by two elements in~$\Tauh$,
gives
\begin{equation} \label{Ione-final}
\begin{split}
\Ione
&\geq \sum_{\K\in \Tauh}(\IKone - |\IKtwo|) \\
&\geq \Big(2\gamma\Ceqtwo^2\Ceqthree^2-\delta\Ceqone^2\Big) \sum_{\K\in \Tauh} \Norm{\epsbfun(\vbfh)}_{\Ltwo(\K)}^2
            - \frac{\Ctilde}{\delta}\sum_{\F\in\Fcalh}\frac{1}{h}{\Norm{\PizF\jump{(\vbfh)_{\tangscalar}}}}_{\Ltwo(\F)}^2.
\end{split}
\end{equation}


We now estimate~$\Itwo$.
Using the identity~\eqref{skmoi-value-on-edge}
and the~$\Ltwo$-orthogonality of~$\PizF$,
we have
\begin{equation} \label{Itwo}
\Itwo 
= \sum_{\F \in \Fcalh} \frac{1}{h}\int_\F \PizF\jump{(\vbfh)_{\tangscalar}} 
  \cdot \jump{(\vbfh)_{\tangscalar}} \ds
= \sum_{\F \in \Fcalh} \frac{1}{h}
  \Norm{\PizF\jump{(\vbfh)_{\tangscalar}}}_{\Ltwo(\F)}^2.
\end{equation}

Combining the estimates~\eqref{Ione-final} and~\eqref{Itwo}
for~$\Ione$ and~$\Itwo$,
gives
\begin{equation*}
\begin{split}
\btwo(\taubfh,&\vbfh)
= \Ione + \Itwo \geq \sum_{\K\in\Tauh}(\IKone - |\IKtwo|) + \Itwo \\
&\geq \Big(2\gamma\Ceqtwo^2\Ceqthree^2-\delta\Ceqone^2\Big)\sum_{\K \in \Tauh}\Norm{\epsbfun(\vbfh)}_{\Ltwo(\K)}^2
  + \Big(1 - \frac{\Ctilde}{\delta}\Big)
  \sum_{\F \in \Fcalh}
  \frac{1}{h}\Norm{\PizF\jump{(\vbfh)_{\tangscalar}}}_{\Ltwo(\F)}^2.
\end{split}
\end{equation*}
By choosing the positive parameters~$\delta$ and~$\gamma$ as
\begin{equation} \label{delta-gamma}
\delta = 2\Ctilde,
\quad \text{and}\quad
\gamma = 2\Ctilde\Big(\frac{\Ceqone}{\Ceqtwo\Ceqthree}\Big)^2,
\end{equation}
and using the identity~\eqref{deveps-equal-eps},
we conclude
\begin{equation*}
\btwo(\taubfh,\vbfh)
\gtrsim  \sum_{\K \in \Tauh} \Norm{\dev(\epsbfun(\vbfh))}_{\Ltwo(\K)}^2
  + \sum_{\F \in \Fcalh}
  \frac{1}{h}\Norm{\PizF\jump{(\vbfh)_{\tangscalar}}}_{\Ltwo(\F)}^2,
\end{equation*}
which yields estimate~\eqref{inf-sup:i}.

We next show~\eqref{inf-sup:ii}.
Using the representation~\eqref{tau-decomposition} of~$\taubfh$
and the orthogonality property~\eqref{orthogonality:dev-sym} twice and
choosing~$\qi = \rkmoi$ in~\eqref{stability-dofs:2},
applying
the estimate~\eqref{skmoi-leq-PizF},
the Cauchy--Schwarz inequality,
the weighted Young inequality with positive parameter~$\eta >0$,
and adding some non-negative terms,
yield
\begin{equation} \label{intermediate-tauhnorm}
\begin{split}
&\Norm{\taubfh}_{\Ltwo(\K)}^2
\overset{\eqref{tau-decomposition}}{\leq}
\sum_{i=1}^3 \Norm{\rkmoi\dev(\normi\otimes\normi)}_{\Ltwo(\Ki)}^2 + \sum_{i=1}^3\Norm{\skmoi\sym(\normi\otimes\tangi)}_{\Ltwo(\Ki)}^2 \\
&= \sum_{i=1}^3 \int_{\Ki}\rkmoi\dev(\normi\otimes\normi):\dev(\normi\otimes\normi)\rkmoi\dx + \sum_{i=1}^3\Norm{\skmoi\sym(\normi\otimes\tangi)}_{\Ltwo(\Ki)}^2 \\
&\overset{\eqref{stability-dofs:2},\eqref{sym:frob-norm}}{=}
\sum_{i=1}^3 - \gamma\int_{\Ki} \epsbfun(\vbfh):\dev(\normi\otimes\normi)\rkmoi \dx
    + \frac12\sum_{i=1}^3 \Norm{\skmoi}_{\Ltwo(\Ki)}^2 \\
&\overset{\eqref{skmoi-value-on-edge}}{\leq}
\sum_{i=1}^3 \gamma\Norm{\epsbfun(\vbfh)}_{\Ltwo(\K)} \Norm{\rkmoi \dev(\normi\otimes\normi)}_{\Ltwo(\Ki)}
    +  \frac{\Ctilde}{2}\sum_{i=1}^3\frac{1}{h} \Norm{\PizFi\jump{(\vbfh)_{\tangscalar}}}^2_{\Ltwo(\Fi)} \\
&{\leq}
\frac{\eta}{2}\gamma^2\Norm{\epsbfun(\vbfh)}^2_{\Ltwo(\K)}
    +  \frac{\Ctilde}{2}\sum_{i=1}^3 \frac{1}{h} \Norm{\PizFi\jump{(\vbfh)_{\tangscalar}}}^2_{\Ltwo(\Fi)} \\
&\qquad    + \sum_{i=1}^3 \frac{1}{2\eta}(\Norm{\rkmoi \dev(\normi\otimes\normi)}^2_{\Ltwo(\Ki)}
        + \Norm{\skmoi \sym(\normi\otimes\tangi)}^2_{\Ltwo(\Ki)}).
\end{split}
\end{equation}
Then,
combining the above display
with
$\dev(\normi\otimes\normi) : \skmoi \sym(\normi\otimes\tangi)= 0$
due to the orthogonality property~\eqref{orthogonality:dev-sym},
we obtain
\begin{equation*}
\begin{split}
\Norm{\taubfh}_{\Ltwo(\K)}^2
&\overset{\eqref{intermediate-tauhnorm},\eqref{orthogonality:dev-sym}}{\leq}
\frac{\eta}{2}\gamma^2\Norm{\epsbfun(\vbfh)}^2_{\Ltwo(\K)}
    +  \frac{\Ctilde}{2}\sum_{i=1}^3 \frac{1}{h} \Norm{\PizFi\jump{(\vbfh)_{\tangscalar}}}^2_{\Ltwo(\Fi)} \\
&\qquad    + \sum_{i=1}^3 \frac{1}{2\eta}\Big(
            \Norm{\rkmoi \dev(\normi\otimes\normi)}^2_{\Ltwo(\Ki)}
          + \Norm{\skmoi \sym(\normi\otimes\tangi)}^2_{\Ltwo(\Ki)} \\
&\qquad   + 2\int_{\Ki} \rkmoi \dev(\normi\otimes\normi) : \skmoi \sym(\normi\otimes\tangi) \dx\Big) \\ 
&= \frac{\eta}{2}\gamma^2\Norm{\epsbfun(\vbfh)}^2_{\Ltwo(\K)}
    + \frac{\Ctilde}{2} \sum_{i=1}^3 \frac{1}{h} \Norm{\PizFi\jump{(\vbfh)_{\tangscalar}}}^2_{\Ltwo(\Fi)}
    + \sum_{i=1}^3 \frac{1}{2\eta}\Norm{\taubfh}^2_{\Ltwo(\Ki)},
\end{split}
\end{equation*}
where the constant~$\gamma$ is fixed as in~\eqref{delta-gamma}.
Using the identity~\eqref{deveps-equal-eps}
in the above display
yields
\begin{equation*} \label{upper-bound-taubfh}
\big(1-\frac{1}{2\eta}\big) \Norm{\taubfh}_{\Ltwo(\K)}^2
\leq \frac12(\eta\gamma^2 + \Ctilde) \Big(\Norm{\dev(\epsbfun(\vbfh))}_{\Ltwo(\K)}^2 + \sum_{i=1}^3 \frac{1}{h}\Norm{\PizFi\jump{(\vbfh)_{\tangscalar}}}^2_{\Ltwo(\Fi)}\Big).
\end{equation*}
Choosing~$\eta = 1$,
both the coefficients
$(1-1/(2\eta))$ and $\eta\gamma^2 + \Ctilde$ are positive;
hence,
summing over the elements~$\K$ in~$\Tauh$,
yields~\eqref{inf-sup:ii}.

The inf-sup estimate~\eqref{discrete-inf-sup:btwo}
follows from~\eqref{inf-sup:i-ii} and
the norm equivalence~\eqref{norm-equivalence-h-dev}.
This concludes the proof.
\end{proof}

\end{document}
